\subsection{Existence of retractions for linear maps}

PROPOSITION 9. Let A be a ring, M a finitely generated A-module, N a projective A-module, and u : M → N an A-linear map.

\begin{enumerate}
\def\labelenumi{\alph{enumi})}
\tightlist
\item
  Let p be a prime ideal of A. The following conditions are equivalent:
\end{enumerate}

\begin{enumerate}
\def\labelenumi{(\roman{enumi})}
\item
  there exists \(f \in \mathrm{A} - \mathfrak{p}\) and \(v \in \operatorname{Hom}_{\mathrm{A}_f}(\mathrm{N}_f, \mathrm{M}_f)\) with \(v \circ u_f = \operatorname{Id}_{\mathrm{M}_f}\) ;
\item
  there exists \(v \in \operatorname{Hom}_{\mathrm{A}_{\mathfrak{p}}}(\mathrm{N}_{\mathfrak{p}}, \mathrm{M}_{\mathfrak{p}})\) with \(v \circ u_{\mathfrak{p}} = \operatorname{Id}_{\mathrm{M}_{\mathfrak{p}}}\);
\item
  the \(\kappa(\mathfrak{p})\text{-linear map }1 \otimes u : \kappa(\mathfrak{p}) \otimes_{\mathrm{A}} \mathrm{M} \to \kappa(\mathfrak{p}) \otimes_{\mathrm{A}} \mathrm{N}\) is injective;
\item
  there exists an integer \(m \geqslant 0\), elements \(x_{1}, \ldots, x_{m}\) of M, and linear forms \(y_{1}, \ldots, y_{m}\) on N such that the images of the \(x_{i}\) in \(\mathrm{M}_{\mathfrak{p}}\) generate the \(\mathrm{A}_{\mathfrak{p}}\)-module \(\mathrm{M}_{\mathfrak{p}}\) and \(\det(<y_{j}, u(x_{i})>) \notin \mathfrak{p}\);
\end{enumerate}

If condition (iv) is satisfied, we have \(m = [\kappa(\mathfrak{p}) \otimes_{\mathrm{A}} \mathrm{M} : \kappa(\mathfrak{p})]\) and the elements \(1 \otimes x_i\) form a basis of the \(\kappa(\mathfrak{p})\)-vector space \(\kappa(\mathfrak{p}) \otimes_{\mathrm{A}} \mathrm{M}\).

\begin{enumerate}
\def\labelenumi{\alph{enumi})}
\setcounter{enumi}{1}
\tightlist
\item
  The set U of prime ideals p of A satisfying the conditions in a) is an open subset of Spec(A), and the following conditions are equivalent:
\end{enumerate}

\begin{enumerate}
\def\labelenumi{(\roman{enumi})}
\item
  We have U = Spec(A);
\item
  U contains all maximal ideals of A;
\item
  there exists \(v \in \operatorname{Hom}_{\mathrm{A}}(\mathrm{N}, \mathrm{M})\) with \(v \circ u = \operatorname{Id}_{\mathrm{M}}\);
\item
  u is injective and Coker(u) is a projective A-module.
\end{enumerate}

Let us prove a).

(i)⇒(ii)⇒(iii): these implications are clear.

(iii)⇒(iv): Set \(m = [\kappa(\mathfrak{p}) \otimes_{\mathrm{A}} \mathrm{M} : \kappa(\mathfrak{p})]\) and let \((x_{1}, \ldots, x_{m})\) be a sequence of elements of M such that the elements \(1 \otimes x_{i}\) form a basis of the \(\kappa(\mathfrak{p})\)-vector space \(\kappa(\mathfrak{p}) \otimes_{\mathrm{A}} \mathrm{M}\). By Nakayama's lemma, the images of the \(x_{i}\) in \(M_{p}\) generate the \(A_{p}\)-module \(M_{p}\). If condition (iii) is satisfied, the elements \(1 \otimes u(x_{i})\) of the \(\kappa(\mathfrak{p})\)-vector space \(\kappa(\mathfrak{p}) \otimes_{\mathrm{A}} N\) are linearly independent.

There exists moreover an A-module \(N'\), an index set I, and an isomorphism of A-modules \(\theta : N \oplus N' \to \mathrm{A}^{(\mathrm{I})}\), from which one deduces an isomorphism of \(\kappa(\mathfrak{p})\)-vector spaces

\[
\overline {{{\theta}}}: (\kappa (\mathfrak {p}) \otimes_ {\mathrm{A}} N) \oplus (\kappa (\mathfrak {p}) \otimes_ {A} N ^ {\prime}) \rightarrow \kappa (\mathfrak {p}) ^ {(I)}.
\]

The elements \(t_i = \bar{\theta}(1 \otimes u(x_i), 0)\) of \(\kappa(\mathfrak{p})^{(I)}\) form a finite free family. There therefore exist elements \(\alpha_1, \ldots, \alpha_{m}\) of I such that one has \(\det(\mathrm{pr}_{\alpha_j}(t_i)) \neq 0\) ; the linear forms \(y_j : z \mapsto \mathrm{pr}_{\alpha_j}(\theta(z, 0))\) on N are suitable.

Suppose condition (iv) is satisfied. Denote \((a_{ij}) \in \mathbf{M}_m(\mathrm{A})\) the matrix of coefficients \(a_{ij} = \langle y_j, u(x_i)\rangle\). Let \(g\) be an element of \(\mathrm{A} - \mathfrak{p}\) such that the images of the \(x_i\) generate the \(\mathrm{A}_g\)-module \(\mathrm{M}_g\) (II, § 5, no. 1, prop. 2), and let \(f = g \det(a_{ij})\). Since \(\det(a_{ij})\) is invertible in \(\mathrm{A}_f\), the images of the elements \(u(x_i)\) in \(\mathrm{N}_f\) are linearly independent; consequently the images of the \(x_i\) in \(\mathrm{M}_f\) form a basis of this \(\mathrm{A}_f\)-module. This proves the last assertion of a). Let us now prove (i). Denote by \(w \in \operatorname{Hom}_{\mathrm{A}}(\mathrm{N}, \mathrm{M})\) the map \(z \mapsto \sum_{j} \langle y_j, z\rangle x_j\). One has

\[
w \circ u (x _ {i}) = \sum_ {j} a _ {i j} x _ {j};
\]

since the images of the \(x_{i}\) form a basis of \(M_{f}\) and the matrix \((a_{ij})\) is invertible in \(\mathbf{M}_{m}(\mathrm{A}_{f})\), the endomorphism \((w \circ u)_{f}\) of \(M_{f}\) is bijective, and the map \(v = (w \circ u)_{f}^{-1} \circ w_{f} \in \operatorname{Hom}_{\mathrm{A}_{f}}(\mathrm{N}_{f}, \mathrm{M}_{f})\) satisfies condition (i).

Let us prove b). The fact that U is open follows from condition (i) of a).

(iii)⇒(i)⇒(ii): this is clear.

(iv)⇒(iii): under the hypotheses of (iv), the sequence 0 → M \(\xrightarrow{u}\) N \(\longrightarrow\) Coker(u) → 0 is exact and split, whence (iii).

\begin{enumerate}
\def\labelenumi{(\roman{enumi})}
\setcounter{enumi}{1}
\tightlist
\item
  (ii) \(\Rightarrow\) (iv): introduce as above an isomorphism of A-modules \(\theta: \mathrm{N} \oplus \mathrm{N}' \to \mathrm{A}^{(\mathrm{I})}\). Let \(u'\) be the map from M to \(\mathrm{A}^{(\mathrm{I})}\) defined by \(u'(x) = \theta(u(x), 0)\). There exists a finite subset J of I such that the image of \(u'\) is contained in the submodule \(\mathrm{A}^{\mathrm{J}}\) of \(\mathrm{A}^{(\mathrm{I})}\). Let \(u'': \mathrm{M} \to \mathrm{A}^{\mathrm{J}}\) be the map induced by \(u'\). Under hypothesis (ii), for every maximal ideal m of A, the \(\mathrm{A}_{\mathfrak{m}}\)-linear map \(u_{\mathfrak{m}}'\) from \(\mathrm{M}_{\mathfrak{m}}\) into \(\mathrm{A}_{\mathfrak{m}}^{(\mathrm{I})}\) admits a retraction, and hence so does \(u_{\mathfrak{m}}''\); thus \(u_{\mathfrak{m}}''\) is injective and its image is a direct summand in \(\mathrm{A}_{\mathfrak{m}}^{\mathrm{J}}\), so that its cokernel is a projective \(\mathrm{A}_{\mathfrak{m}}\)-module. The A-module \(\operatorname{Coker}(u'')\) is of finite presentation by construction; it is therefore projective (II, § 5, no. 2, th. 1). The homomorphism \(u''\) is injective (II, § 3, no. 3, th. 1); consequently, \(u\) is injective. The A-module \(\operatorname{Coker}(u')\) is isomorphic, on the one hand, to \(\operatorname{Coker}(u) \oplus \mathrm{N}'\) and, on the other hand, to \(\operatorname{Coker}(u'') \oplus \mathrm{A}^{(\mathrm{I}-\mathrm{J})}\). Since the A-modules \(\mathrm{A}^{(\mathrm{I}-\mathrm{J})}\), \(\operatorname{Coker}(u'')\), and \(\mathrm{N}'\) are projective, it follows that \(\operatorname{Coker}(u)\) is projective, which completes the proof of (iv).
\end{enumerate}
% semantic-fragments: legacy-input-seam
\relax{}
